\documentclass[compress]{beamer}
\useinnertheme{rectangles}
\usecolortheme{orchid}
\useoutertheme{miniframes}
\setbeamertemplate{blocks}[rounded][shadow=true]
\setbeamercolor{separation line}{use=structure,bg= structure.fg!20!bg}
\setbeamercovered{dynamic}
\newcommand{\lined}{\hfill\hrule\hfill\vspace{1.5mm}}
\setbeamertemplate{headline}{%
  \textcolor{blue}{\lined}
  \textcolor{blue}{\lined}
  \textcolor{blue}{\lined}
}

\usepackage[english]{babel}
\usepackage[T1]{fontenc}
\usepackage[utf8]{inputenc}
\usepackage{multirow}
\usepackage{multicol}
\usepackage{graphicx}
\usepackage{bm}
\usepackage{color}
\usepackage{amsmath,amssymb,mathtools}

\setbeamertemplate{theorems}[numbered]
\setbeamertemplate{footline}[page number]

\newcommand{\argmin}{\ensuremath{\operatornamewithlimits{argmin}}}
\newcommand{\tR}{\mathbb{R}}
\newcommand{\Frechet}{Fr\'echet}
\newcommand{\Zbold}{\bm{Z}}
\newcommand{\Xbold}{\bm{X}}

\begin{document}

%========================= TITLE =========================
\begin{frame}

\title{Discussion of \\ \emph{Deep \Frechet{} Regression} \\ \vspace{0.15cm} {\small Su I Iao, Yidong Zhou \& Hans-Georg M\"uller}}

\author{Edward Gunning, Giles Hooker, Kyle Stanley \& Jeffrey S. Morris}

\date{JASA Special Issue on Statistical Science in Artificial Intelligence \\ \vspace{0.1cm} {\small August 2026}}

\maketitle

\end{frame}


%========================= SLIDE 1 =========================
\begin{frame}{A Program for Data in Metric Spaces}

\begin{itemize}
\item Modern responses are complex objects: distributions, covariance matrices, graph Laplacians, trees, shapes; do not fit within current statistical paradigms. 

\item Papers continue influential program on \Frechet{} regression for random objects (Petersen \& M\"uller 2019; Zhou \& M\"uller 2022; Dubey, Chen \& M\"uller 2024).

\item Assume {\em only} a distance $d(\cdot,\cdot)$ between objects. Two strategies
\begin{itemize}
\item[(a)] {\bf Weight the distances} --- \Frechet{} means / regression
\[
\hat{m}(x) = \argmin_{y \in \Omega} \sum_i w_i(x)\, d^2(Y_i,y)
\]

\item[(b)] {\bf Translate to Euclidean space}, model there, map back.
\end{itemize}
\end{itemize}

Use deep learning for both; we will focus on (b) composite map $m = v \circ g_0$, a useful mirror to our own work on latent representations.

\end{frame}


%========================= SLIDE 2 =========================
\begin{frame}{How Do We Choose the Manifold Dimension?}

Embed the responses, then predict the embedding with a deep network:
\[
\underbrace{Y \xrightarrow{\ \psi \ (\text{ISOMAP})\ } \Zbold = (Z_1,\dots,Z_r) \in \tR^r}_{\text{embed responses}},
\qquad
\hat{Z}_j = \hat{g}_j(\Xbold).
\]

\begin{itemize}
\item Paper and supplement: $r = 2$; accompanying code fixes $r \le 2$.

\item Exact in simulations (Gaussian $\Rightarrow$ 2-D) but unknown in practice, and sensitivity  unexplored.
\end{itemize}

\end{frame}


\begin{frame}{Compact Near-Lossless Latent Representations} 

{\bf CLaRe} (Zohner, Gunning, H.\ \& Morris 2025) examines how to choose a Euclidean representation $f(y)$

\vspace{0.2cm}

Assume encoder $f(y): \Omega \rightarrow \mathbb{R}^r$ and decoder $g(x): \mathbb{R}^r \rightarrow \Omega$, chosen according to

\begin{itemize}
\item {\bf Near-losslessness}: object reconstruction, controlled at the {\em tail}:
\[
\mathbb{P}\{\, d(\,g(f(Y)),\,Y\,) \le \epsilon \,\} \ge 1-\alpha.
\]

\item {\bf Compactness}: as few coordinates as possible.
\end{itemize}

Reconstruction error measured out-of-sample, or in cross-validation. 

\vspace{0.2cm}

Reasoning: represent all object details, move analysis to multivariate, Euclidean settings. 

\vspace{0.2cm}

Metric space analogy: examine reconstruction of distances. 

\end{frame}


%========================= SLIDE 3a =========================
\begin{frame}{A Variation on Simulation Studies}

Author's simulated quantile functions; use same conditional mean and variance, add a {\em skew} parameter $\alpha$ determined by $X_9$:

\begin{itemize}
\item Raw quantile functions of the three settings look essentially alike.
\item Project onto $(1,\Phi,\Phi^2) \approx (\text{mean, variance, skew})$ color by $X_9$: effect of skew normal becomes clear.  
\end{itemize}

\centering
\includegraphics[width=0.96\textwidth]{figures/manifold-data-group-third-isomap.png}

\end{frame}


%========================= SLIDE 3b =========================
\begin{frame}{The Signal in the Tail}

Reconstruct each object's geodesic distances from an $r$-dimensional ISOMAP embedding; look at the {\em individual} errors.

\centering
\includegraphics[width=\textwidth]{figures/bottom-row-plot.pdf}

{\raggedright
\begin{itemize}
\item Aggregate summaries at $r = 2$ look ok. 
\item {\em tail} collapses at $r = 3$; but this is what captures skew. 
\end{itemize}
\par}

\end{frame}


%========================= SLIDE 4a =========================
\begin{frame}{Ups and Downs of Generality}

%\begin{center}\bf Where does the error live?\end{center}

Multi-stage process produces errors in two places
\begin{itemize}
\item ISOmap embedding/reconstruction

\item Predictive error from DNN 
\end{itemize}
Appears in mathematical analysis, but also model checking, uncertainty quantification, $\ldots$.. 

\vspace{0.2cm}

\begin{itemize}
\item CLaRe objective:  make embedding/reconstruction error ignorable, all error in Euclidean representation
\begin{itemize}
\item Requires encoder/decoder architecture; not obviously implied with just metric space assumptions

\item Reconstruction of {\em distances} does not imply reconstruction of {\em objects}
\end{itemize}

\item Zhou, Lao,  Muller, 2026: parameterize similarity weights $w_j(X)$, $\Rightarrow$ error in object space
\begin{itemize}
\item Does not require decoder, less obvious model validation. 
\end{itemize}
\end{itemize}


%DFR's generality spreads error across {\em stages}:
%\[
%d\{\hat{m}(\Xbold),m(\Xbold)\} = O_p\!\Big(
%\underbrace{h^{\frac{2}{\gamma_1-1}}}_{\text{LFR bias}}
%+ \underbrace{(nh^r)^{-\frac{1}{2(\gamma_2-1)}}}_{\text{LFR variance}}
%+ \underbrace{(h^{-r-1}\zeta_n)^{\frac{1}{\gamma_3-1}}}_{\text{embed + DNN}}
%\Big)
%\]
%with $\zeta_n = \kappa_n \log^{3/2} n + u_n$; the ISOMAP bias $u_n$ has, as the authors note, no asymptotic theory yet.

%\begin{itemize}
%\item Error sits in {\em both} the latent prediction and the \Frechet{} decode $\Rightarrow$ hard to attribute, hard to do UQ.
%
%\item CLaRe's aim: drive the reconstruction error to (near) zero, so {\em all} uncertainty is the Euclidean latent model.
%
%\item Good {\em distance} reconstruction $\ne$ good {\em object} reconstruction --- in a bare metric space only the former is even defined.
%\end{itemize}

\end{frame}


%========================= SLIDE 4b =========================
\begin{frame}{Does Generality Preclude Decoders}

\Frechet{} regression defined by optimization {\em over objects}:
\[
\hat{m}(\Xbold) = \hat{v}_h(\hat{\Zbold})
= \argmin_{y \in \Omega} \sum_{i=1}^n \hat{w}(\hat{\Zbold}_i,\hat{\Zbold},h)\, d^2(Y_i,y).
\]

\begin{itemize}
\item To {\em search} over $y \in \Omega$ you must already represent objects on a computer (discretized quantile functions, PSD parametrizations, \dots).

\item Mathematical generality does not circumvent object-specific implementation; user must specify how to optimized within $\Omega$

\item Computer representation means explicit encoder--decoder is available in principle. 

\item Once the encoder/decoder is (near-)lossless, we are back in familiar, Euclidean territory. 
\end{itemize}


But CLaRe does not admit the mathematical generality that simply working with metric spaces allows.  

\end{frame}


\begin{frame}{References}
    
\end{frame}

\end{document}
